\documentclass[11pt,a4paper]{article}

\usepackage[T1]{fontenc}
\usepackage[utf8]{inputenc}
\usepackage[english]{babel}
\usepackage{lmodern}
\usepackage{amsmath,amssymb}
\usepackage{graphicx}
\usepackage{booktabs}
\usepackage{microtype}
\usepackage[margin=2.4cm]{geometry}
\usepackage[hidelinks]{hyperref}
\usepackage{natbib}
\bibliographystyle{plainnat}

\graphicspath{{figures/}}
\title{\textbf{Measuring the Diversity a Feed Actually Exposes}\\[2mm]
       \large A standard, its attacks, and what it takes to verify it}
\author{Sylvain Geffroy}
\date{August 2026 --- \emph{working paper, open to review}}

\begin{document}
\maketitle

\begin{abstract}
\noindent
A regulator seeking to constrain a recommender algorithm needs a quantity it can observe
without access to the platform's code. This note proposes one --- the \textbf{Exposed
Diversity Index} (EDI), the entropy of the items served, projected onto a declared catalogue
of viewpoints and weighted by the attention each rank commands --- and then subjects it to the
attacks it must withstand to serve as a standard.

Three of those attacks succeed against the naive formulations, and each forces a correction. A
floor bearing on a feed's \emph{labels} is saturated at zero cost: a platform able to decouple
label from content obtains full marks for strictly zero content diversity. Rao's quadratic
entropy, the first replacement considered, has a \emph{bimodal} constrained optimum: it would
prescribe the very polarisation it claims to measure. And any rank-blind measure is satisfied
by \emph{burying} the divergent items --- a platform certified at $0.70$ exposes only $0.36$ of
real diversity, and closing the loophole doubles the engagement cost.

Verifying such a standard on real data requires knowing \emph{exposure}. We give an exact
within-feed test that decides whether a recommender log permits any exposure correction at
all, and an estimator of position-bias severity. Applied to three public logs, they establish
that MIND does not preserve the served order --- the estimator returns five incompatible
severities there --- while Baidu-ULTR does ($\hat\eta = 1.10 \pm 0.09$), and the Open Bandit
Dataset permits this work's only confrontation with a ground truth: the value of a
never-deployed policy is recovered to within $2.5\%$, against $31.6\%$ error for the naive
estimate.

No public dataset carries both the served rank and an interpretable viewpoint label. We
therefore specify, under Article~40 of the \emph{Digital Services Act}, four aggregate tables
containing no personal data, verified to recompute the measurements identically for a hundred
times fewer rows than a raw log.

Part~\ref{part:model} presents the thermodynamic model the index came from, and
\S\ref{sec:analogy} reports what remains of it: the quantum decoherence analogy that started
this work survives none of its transfers. All code, tests and figures are reproducible in
containers.
\end{abstract}

\part{The measurement}
\label{part:measure}

\section{The metrological problem}
\label{sec:problem}

Regulating a recommender algorithm presupposes an observable quantity. Counting problematic
items will not do: the indicator lags, it is gameable, and it requires ruling on the substance
of every item. The quantity adopted here is of a different nature --- it bears on the
\emph{shape} of exposure, not on the content of messages:

\begin{quote}
What share of the attention served to a reader goes to each of the viewpoints the platform
could have served?
\end{quote}

The question has three virtues: it can be measured from outside, it requires no judgement on
the truth of any item, and it admits a threshold expressed as a percentage, hence transferable
across platforms.

It also has three traps, and this note is organised around them.

\begin{enumerate}
\item \textbf{A label is not an item.} Measuring the diversity of \emph{announced categories}
      leaves the platform to choose what is announced.
\item \textbf{A composition is not an exposure.} A feed may contain the required diversity and
      place it where nobody looks.
\item \textbf{A logged click is not a preference.} Observed clicks were produced by the
      platform's ranking, not by the one under evaluation.
\end{enumerate}

Each trap is measurable, and each was measured. What survives all three is the retained
definition.

\section{The exposed diversity index}
\label{sec:index}

Given a feed of $n$ positions, a reference catalogue of $k$ viewpoints declared by the
regulator, and an attention discount $w_R$ attached to rank $R$, the \emph{exposed}
distribution is
\begin{equation}
q_i = \frac{\sum_{R=1}^{n} w_R \, \mathbb{1}\!\left[\text{the item served at rank } R
      \text{ falls in bin } i\right]}{\sum_{R=1}^{n} w_R},
\label{eq:distribution}
\end{equation}
and the index is its normalised Shannon entropy:
\begin{equation}
\mathrm{EDI} = \frac{H(q)}{\log_2 k}, \qquad H(q) = -\sum_{i=1}^{k} q_i \log_2 q_i .
\label{eq:edi}
\end{equation}

Three choices distinguish \eqref{eq:distribution} from the entropy one would write
spontaneously, and each follows from an attack that succeeded (\S\ref{sec:adversarial}).

\begin{center}
\begin{tabular}{@{}p{0.30\linewidth}p{0.30\linewidth}p{0.32\linewidth}@{}}
\toprule
\textbf{Choice} & \textbf{Reason} & \textbf{Without it} \\
\midrule
$q$ bears on the items served, projected onto the catalogue's bins & a label is chosen, an item
is observed & the index reaches $1.000$ for zero content diversity, at zero cost \\[2mm]
each rank is weighted by $w_R$ & a reader consults the first item far more often than the last
& the standard is satisfied by burial: $0.70$ certified, $0.36$ exposed \\[2mm]
the denominator $\log_2 k$ comes from the declared catalogue & comparing platforms requires the
same unit & a closed feed shows one modality, the denominator degenerates and the index
flatters \\
\bottomrule
\end{tabular}
\end{center}

\paragraph{Attention discount.} The results below use $w_R = 1/R$, the reciprocal-rank
discount standard in ranking evaluation. The choice of $w$ belongs to the regulator, as does
the catalogue; what matters here is that it be \emph{declared}, since it is what makes the
standard non-circumventable by order.

\paragraph{Associated control quantity.} The gap between \eqref{eq:edi} computed on the feed's
composition ($w_R$ constant) and on its exposed distribution ($w_R = 1/R$) is zero for a
platform that does not relegate, and grows with burial. It is the only quantity in this work
that compares a measure \emph{with itself}, and hence the only one that thresholds without an
external reference.

\paragraph{How to publish the figure.} A normalised entropy is not a diversity: it is not linear
in what we mean by "twice as diverse". Its conversion into an \textbf{effective number of
viewpoints}, ${}^{1}\!D = k^{\mathrm{EDI}}$, is \citep{jost2006}: a floor of $0.70$ on a catalogue
of four is $2.6$ effective viewpoints, and the $0.44$ actually exposed by the burying feed is
$1.9$.

\paragraph{Protocol.} The unit is the item \emph{served}, not the item consulted: the platform
is audited, not the reader. The index is computed per reader over a sliding window and then
aggregated; the regulatory quantity is the \emph{share of the population below the threshold},
not the mean, since a satisfactory mean can hide an entirely enclosed minority.

\section{The adversarial test}
\label{sec:adversarial}

A standard is judged by what it forces a rational adversary to do. We therefore place the
platform in the role it will occupy: maximising engagement \emph{subject to} the constraint,
rather than complying in good faith.

\paragraph{Method.} A feed has $n$ positions to fill from a catalogue of $k$ viewpoints, hence
$k^n$ possible feeds. For $n = 8$ and $k = 4$, all $65{,}536$ feeds are \emph{enumerated}: the
optimum obtained is exact, not the output of a heuristic. This is not a luxury. These are
negative results --- standards that fail --- and an optimum missed by a solver would look
exactly like a standard that holds.

\subsection{A label-based standard saturates at zero cost}

If the platform controls the label independently of the item, the entropy of labels stops
constraining the entropy of items. Optimisation confirms it without nuance: under full
decoupling the platform obtains an index of $1.000$ --- full marks --- for strictly zero content
diversity, without giving up a point of engagement. One need not go that far: at half
decoupling, the constraint retains only $36\%$ of its force.

The consequence is direct: \textbf{the measure must bear on the items served}. It calls for a
complementary diagnostic, the \emph{signature excess} --- the gap between two indices computed
on the same feed, relative to what an honest catalogue would show at the same index --- which is
zero for an honest platform and grows with decoupling.

\subsection{The first replacement prescribed polarisation}

Rao's quadratic entropy \citep{rao1982} --- the mean distance between two items drawn from the
feed, known in recommendation as \emph{intra-list distance} --- appeared to answer the
objection: it sees items, not labels. Yet its engagement-constrained optimum is
\textbf{bimodal}: at a given required diversity, the cheapest way to supply it is to serve two
extreme blocks and empty the centre. A standard built on it would prescribe the polarisation it
claims to measure. The defect is known in the diversification literature
\citep{ohsaka2023}; it is recovered here as a direct consequence of constrained optimisation.

The retained floor therefore bears on an \emph{entropy}, computed on the items served projected
onto the catalogue's bins --- a measure whose constrained optimum is spread rather than bimodal.

\subsection{A rank-blind standard is circumvented by burial}

The measures above bear on the feed's \emph{composition}. Retested on \emph{ordered} feeds, all
four candidate diversity measures are circumvented the same way: serve the divergent items at
the last ranks, where attention no longer goes.

\begin{center}
\begin{tabular}{@{}lrrrr@{}}
\toprule
\textbf{Measure} & \textbf{Cost} & \textbf{Displayed} & \textbf{Exposed} & \textbf{Cost if rank-aware} \\
\midrule
Rao's entropy & $8.2\%$ & $0.750$ & $\mathbf{0.355}$ & $18.9\%$ \\
Position entropy & $10.7\%$ & $0.774$ & $\mathbf{0.443}$ & $20.9\%$ \\
Gaussian ILD & \multicolumn{4}{l}{floor unattainable} \\
Target proximity & $5.9\%$ & $0.750$ & $\mathbf{0.628}$ & $10.6\%$ \\
\bottomrule
\end{tabular}
\end{center}

The optimal feed always has the same shape: items of the preferred viewpoint on top, divergent
ones relegated. \textbf{A platform certified at $0.70$ exposes only $0.36$.} A rank-aware floor
closes the loophole and \emph{doubles} the engagement cost: a standard that cost no more would
close nothing.

The gap grows with the requirement --- $0.262$ at a floor of $0.40$, $0.395$ at $0.70$ for Rao's
entropy: the more diversity a rank-blind standard demands, the more profitable burial becomes.

\begin{figure}[t]
\centering
\includegraphics[width=\linewidth]{fig15_rang_adverse.png}
\caption{Displayed versus exposed diversity under a rank-blind floor; the gap by floor level;
position-bias severity recovered by fixed effects, and the exploration threshold below which it
ceases to be identifiable; the cost of a posited $\eta$.}
\label{fig:adversarial}
\end{figure}

\section{Measuring exposure}
\label{sec:exposure}

Everything above presupposes knowing what attention goes to which rank. The standard model
posits an examination probability decreasing with rank \citep{joachims2017},
\begin{equation}
e(R) = R^{-\eta},
\label{eq:examination}
\end{equation}
whose severity $\eta$ was, in the first version of this work, \emph{posited}. It can be
estimated.

\subsection{Estimation by item fixed effects}

Under \eqref{eq:examination}, click probability factorises as $P(\text{click} \mid i, R) =
g(i)\, R^{-\eta}$, hence
\begin{equation}
\log \mathrm{CTR}(i, R) = \log g(i) - \eta \log R .
\label{eq:fixedeffects}
\end{equation}
Relevance $g(i)$ is an item fixed effect: it is not estimated but eliminated by within-item
centring. What remains is the only variation that identifies $\eta$ --- that of the \emph{same
item seen at different ranks}. This is the simplest form of intervention harvesting
\citep{agarwal2019}: it requires no experiment, only that the platform did not always rank the
same items in the same places.

On simulated data the estimator recovers $\eta$ from $0.4$ to $1.6$. Under a deterministic
policy it \emph{refuses} to answer rather than invent a figure: no item changes rank, and the
severity is not in the data.

Positing $\eta$ instead of estimating it is not innocuous: where the true cost of a diversity
filter is $6.6\%$, a posited $\eta$ of $0.5$ understates it by $59\%$ and a posited $\eta$ of
$2.0$ overstates it by $179\%$ --- the very order of magnitude of the bias one claimed to
correct.

\subsection{The check that must come first: exchangeability}

The refusal described above catches the visible case --- no rank variation --- and misses the
invisible one: \emph{abundant and artificial} variation. What is needed is a check that counts
not the variation but the \emph{informativeness} of the order.

For a feed of length $L$ carrying $k$ clicks, let $u_R = (R - 1/2)/L$ be the normalised rank.
Under the hypothesis that clicks are indifferent to position \emph{given the feed}, the clicked
positions are a draw without replacement among the $L$ positions, with exactly known expectation
and variance:
\begin{equation}
\mathbb{E}\Big[\sum_{\text{clicks}} u_R\Big] = k\,\bar{u}, \qquad
\mathbb{V}\Big[\sum_{\text{clicks}} u_R\Big] = \frac{k(L-k)}{L-1}\,\sigma^2_u .
\label{eq:exchangeability}
\end{equation}
Conditioning on the feed immunises the test against the mixture of feed lengths, the mean
quality of a feed's items, and its reader's appetite for clicking. Position bias concentrates
clicks at the top: it makes the standardised deviation \emph{negative}, so the sign of the
rejection is itself a verification.

A test that rejects nothing says nothing until one knows what it would reject. Calibration is
done by simulation at identical feed structure: on MIND's, the test detects $\eta = 0.02$ at
twelve standard deviations, and the minimum detectable severity is $\eta \approx 0.004$.

\section{Evaluating a re-ranking}
\label{sec:evaluation}

Measuring a feed is one thing; estimating what a different ranking would cost is another. The
available clicks were produced by the platform's policy, not by the one under evaluation.

\paragraph{Replay is wrong, by a considerable margin.} Replaying logged clicks on a re-ranked
feed is off by $201\%$ at the median, up to $851\%$, and the sign of the error is not
guaranteed: across sixty draws at identical configuration it overstates in fifty-six cases and
understates in four. A naive figure is therefore not even an upper bound.

\paragraph{The correction and its price.} Inverse propensity weighting is unbiased provided the
logging policy is known and gives every action the target policy might choose a chance. Its
price is variance, hence the usual variants --- self-normalisation \citep{swaminathan2015},
weight clipping --- and the diagnostic that must accompany any figure: the \textbf{effective
sample size}, which states how many observations actually carry the estimate.

\paragraph{A confrontation with ground truth.} The Open Bandit Dataset \citep{saito2020}
contains two buckets served in parallel by two different policies, and publishes the true
propensities. One can therefore estimate the value of the uniform policy \emph{from another
policy's data alone}, and compare it with its value measured where it actually served.

\begin{center}
\begin{tabular}{@{}lrr@{}}
\toprule
\textbf{Estimator} & \textbf{Value} & \textbf{Error} \\
\midrule
Ground truth (measured, $452{,}949$ impressions) & $0.005124$ & --- \\
Naive (observed click rate) & $0.006743$ & $+31.6\%$ \\
IPS & $0.005253$ & $\mathbf{+2.5\%}$ \\
Self-normalised & $0.004768$ & $-7.0\%$ \\
IPS clipped at $10$ & $0.004473$ & $-12.7\%$ \\
\bottomrule
\end{tabular}
\end{center}

The correction works. The diagnostic forbids celebrating it: the effective sample size is
$1{,}513$ for $4{,}077{,}727$ impressions, i.e. $0.04\%$. The estimate is unbiased and rests on
the equivalent of fifteen hundred observations; that the error lands at $2.5\%$ owes as much to
luck as to method. This is the real case that justifies the rule: publish the effective size
next to the figure.

\section{Three public logs}
\label{sec:logs}

The instruments of \S\ref{sec:exposure} and \S\ref{sec:evaluation} decide whether a dataset
permits the measurement. Applied to three public logs, they give three different answers.

\subsection{MIND: an order that is not one}

The reference dataset for news recommendation \citep{wu2020} shows ideal rank coverage: a median
item appears at sixteen distinct positions. Yet test \eqref{eq:exchangeability} detects nothing
--- $z = +0.12$ ($p = 0.91$) across $156{,}965$ feeds, replicated at $z = +0.28$ on the second
split --- in a dataset where it would detect $\eta = 0.02$ at twelve standard deviations. The
recorded order is indistinguishable from a shuffle.

The click-rate curve by position nonetheless falls from $0.108$ to $0.038$ and fits
$\hat\eta = 0.39$: a \emph{composition artefact}, deep positions existing only in long feeds,
whose click rate per item is mechanically lower. At fixed feed length the slope is zero and
changes sign from one length to the next.

Finally, estimator \eqref{eq:fixedeffects} accepts the dataset, declares itself identifiable,
and returns a severity that is an increasing function of a mere nuisance threshold: $-0.131$ at
five impressions, $+0.053$ at twenty, $+0.255$ at a hundred, each with a standard error below
$0.007$. \textbf{Five confident and incompatible figures from one dataset.}

Shuffling does not debias the clicks: they were produced under the real order and carry its bias
in full. What the shuffle removes is the \emph{rank}, the only regressor that would allow
correction. On a simulated log of true severity $1.00$, the estimate moves from $1.003$ to
$-0.003$ after shuffling, at identical clicks; evaluating a re-ranking becomes vacuous by
construction.

\subsection{Baidu-ULTR: the positive control}

On a slice of $524{,}164$ documents served across $64{,}200$ search sessions \citep{zou2022},
the same test rejects at $z = -205.7$ ($p < 10^{-12}$) and on the side theory prescribes.
Severity is $\hat\eta = 1.10 \pm 0.09$ by document fixed effects, against $1.49$ by aggregate
fit --- the gap being the quality confounder, the platform placing the best documents on top.
Coverage remains thin: $55$ documents carry the estimate, out of $444{,}709$, because the same
URL rarely reappears.

\subsection{Open Bandit Dataset: position measured without a model}

The random bucket allocates items to positions by chance: the click-rate difference between
positions is a \emph{causal} effect, with no examination model to posit. On a three-thumbnail
horizontal banner it is $\hat\eta = 0.037$ ("all" campaign, $1.37$ million impressions) to
$\hat\eta = 0.105$ ("men" campaign).

\textbf{$\eta$ is therefore a property of the surface, not a constant}: an order of magnitude
separates a vertical results page from a three-thumbnail banner. Transporting one to the other
is exactly the error whose cost \S\ref{sec:exposure} quantifies.

\begin{figure}[t]
\centering
\includegraphics[width=\linewidth]{fig17_rang_servi.png}
\caption{What a genuinely served rank does to the click-rate curve; the exchangeability test on
MIND and on Baidu-ULTR; the position effect measured without a model on the Open Bandit
Dataset's random bucket; the counterfactual estimator judged against the value it estimates.}
\label{fig:logs}
\end{figure}

\subsection{What these three logs do not permit}

None carries both the \emph{served rank} and an \emph{interpretable viewpoint label}. MIND has
editorial categories without the rank; Baidu-ULTR has the rank without labels; the Open Bandit
Dataset has rank, propensity and three categorical attributes, but anonymised --- and a diversity
computed over hashes says nothing, the reference catalogue being a political declaration.

The index of \S\ref{sec:index} has therefore never been measured on a real feed. This is not a
gap in method: it is a gap in data.

\section{What to ask for}
\label{sec:request}

Article~40 of the \emph{Digital Services Act} \citep{dsa2022}, made operational by Delegated
Regulation (EU) 2025/2050 \citep{delegue2025}, opens platform data to vetted researchers. A
request is admissible only if it is \emph{necessary and proportionate}. "Give us your logs" is
neither: it demands personal data the analysis has no need of, and it hands over the easiest
ground for refusal --- service security and trade secrets.

We therefore specify the request as \textbf{four aggregate tables}:

\begin{enumerate}
\item \textbf{feed profiles} --- for each served feed shape (the list of ranks it occupies) and
      each click count, the number of feeds;
\item \textbf{clicks by rank} --- for each profile, each feed click count and each rank, the
      number of clicks;
\item \textbf{(item, rank) cells} --- impressions and clicks, and the service propensity if the
      platform knows it;
\item \textbf{exposure by (rank, viewpoint)} --- impressions and clicks, over the catalogue
      declared by the regulator.
\end{enumerate}

No row designates a reader: a feed appears only by its shape and by the number of clicks it
received.

\paragraph{Sufficiency verified, not asserted.} Test \eqref{eq:exchangeability} is recomputed
from tables 1 and 2 alone to within $3 \times 10^{-12}$ on MIND as on Baidu-ULTR; severity
\eqref{eq:fixedeffects} is recomputed \emph{exactly} from table 3; both diversity measures from
table 4, which therefore suffices to observe burial --- at identical composition, a feed
compliant with the rank-blind floor exposes only $0.160$ of attention to minority viewpoints,
against $0.375$ in composition.

\paragraph{Proportionality quantified.} What is asked weighs $5{,}543$ rows for the $524{,}164$
documents served in the Baidu-ULTR slice (ratio $95$) and $57{,}906$ rows for MIND's
$5{,}843{,}444$ served items (ratio $101$).

\paragraph{The confidentiality threshold is not neutral.} Suppressing low-count cells protects
readers and moves the estimates: going from $5$ to $20$ impressions per cell moves the estimated
severity from $1.10$ to $1.40$, i.e. $+27\%$, because rare cells are those of deep ranks. The
threshold must be fixed in the request and published with the result.

\paragraph{Two simpler aggregation keys were discarded} because they were wrong without being
visible: indexing feeds by length rather than by rank profile --- wrong as soon as a surface skips
ranks --- and omitting the feed click count, which prevents excluding fully clicked feeds and
shifts the standardised deviation from $-205.7$ to $-203.9$. Both produced figures of the right
order of magnitude.

\begin{figure}[t]
\centering
\includegraphics[width=\linewidth]{fig18_demande_article_40.png}
\caption{What the request weighs against the log; the gap between direct and aggregate
measurement; the effect of the confidentiality threshold on estimated severity; burial observed
from the exposure table alone.}
\label{fig:request}
\end{figure}

\section{A filter that optimises the index}
\label{sec:filter}

A classical recommender orders items by estimated relevance. The filter proposed here adds a
term:
\begin{equation}
S(i, c) = \text{Relevance}(i, c) + \mu \, \Delta H(i, c), \qquad \mu \geq 0,
\label{eq:score}
\end{equation}
where $\Delta H$ is the change in index item $i$ would bring to feed $c$. At $\mu = 0$ the
ranking is \emph{identical} to an engagement filter's: the transformation is a parameterised
addition, not a redesign. The parameter $\mu$ follows an annealing schedule --- rising when the
index falls below the threshold, falling once it is restored --- so that intervention is
intermittent.

\paragraph{Status.} This filter is validated against its own model, with synthetic relevance and
four viewpoints. Its real cost in perceived relevance is not measured, and cannot be on public
data as things stand (\S\ref{sec:logs}). What bounds the discussion is the figure of
\S\ref{sec:adversarial}: a rank-aware floor costs $10.6$ to $20.9\%$ of engagement depending on
the measure, on eight-slot feeds enumerated exhaustively.

\subsection{What it is worth, compared with baselines}
\label{sec:baselines}

A filter is not judged against relevance ranking, from which it differs by construction, but
against what a practitioner would do without theory. We therefore compare it with four baselines
--- serving viewpoints in turn, \emph{maximal marginal relevance} \citep{carbonell1998}, a
Boltzmann draw on relevance, and random selection --- and above all with the \textbf{exact
frontier} of the (exposed diversity, engagement) plane, obtained by enumerating the $151{,}200$
ordered arrangements of six items drawn from a pool of ten.

Across $150$ pools, at a fixed floor, the median shortfall is as follows.

\begin{center}
\begin{tabular}{@{}lrrrrr@{}}
\toprule
\textbf{Method} & $0.50$ & $0.60$ & $0.70$ & $0.80$ & $0.90$ \\
\midrule
Round robin & $0.0\%$ & $0.0\%$ & $0.0\%$ & $3.8\%$ & $4.8\%$ \\
MMR & $0.0\%$ & $0.0\%$ & $0.0\%$ & $\mathbf{0.0\%}$ & $\mathbf{0.2\%}$ \\
Boltzmann & $0.0\%$ & $0.0\%$ & $0.1\%$ & $1.2\%$ & $3.7\%$ \\
Filter \eqref{eq:score} & $0.0\%$ & $0.0\%$ & $0.0\%$ & $\mathbf{0.0\%}$ & $\mathbf{1.0\%}$ \\
Random & $9.3\%$ & $9.8\%$ & $11.0\%$ & $12.9\%$ & $19.1\%$ \\
\bottomrule
\end{tabular}
\end{center}

The filter therefore holds the frontier --- which was not a given, random selection reaching the
same diversities while leaving a fifth of the engagement on the table --- but \textbf{it is not
alone there}. A heuristic published in 1998 holds it too, and beats the filter head to head at
floor $0.80$. The filter only stands out at the high requirement: at floor $0.90$, MMR finds no
compliant setting in $47\%$ of pools, against $12\%$ for the filter, because the latter optimises
the constrained quantity where MMR optimises a proxy that saturates.

\paragraph{The price of the standard depends on the reader.} The comparison above draws relevance
independently of viewpoint. Varying that \emph{alignment}, the cost of the $0.90$ floor moves from
$3.8\%$ of engagement --- a reader whose interests cut across viewpoints --- to $17.1\%$ --- a
reader whose preference \emph{is} a viewpoint. The standard therefore costs most exactly where it
serves most, which predicts where the objection will fall. This dependence reconciles the figures
of \S\ref{sec:adversarial}, where relevance was attached to the viewpoint, with those of this
paragraph.


\part{The background model}
\label{part:model}

What precedes stands without what follows. The index, its attacks, the exchangeability test and
the counterfactual estimators depend on no hypothesis about opinion dynamics: they are
measurements on served feeds and observed clicks.

This part presents the model the index \emph{came from}. It has two functions, and only one is
scientific. The first is to account for provenance: the quantity adopted in \S\ref{sec:index}
was not picked at random among possible diversity measures, it comes from a formalism that
describes opinion as a system with a temperature. The second is to state what that formalism is
worth once measured --- and \S\ref{sec:analogy} concludes that the analogy which suggested it
survives none of its transfers.

No result in Part~\ref{part:measure} rests on this part.

\section{Entropic formalism}

\subsection{Two entropies, one quantity}

The von Neumann entropy \citep{vonneumann1932}
\begin{equation}
S(\rho) = -\operatorname{Tr}(\rho \ln \rho)
\end{equation}
vanishes for a pure state. It is computed by diagonalisation: the eigenvalues of $\rho$
form a classical distribution to which the Shannon entropy \citep{shannon1948} applies,
\begin{equation}
H(X) = -\sum_i p_i \log_2 p_i .
\end{equation}
This is precisely what licenses the bridge between the two disciplines: von Neumann entropy
\emph{is} the Shannon entropy of the mixture revealed in the eigenbasis.

\subsection{A necessary clarification}

One common and incorrect formulation must be set aside at once. The evolution of a
\emph{closed} quantum system is unitary, so its von Neumann entropy is constant. What grows
under decoherence is the entropy of the \emph{reduced subsystem}, obtained by tracing over
environmental degrees of freedom.

A clarification follows that strengthens the analogy rather than weakening it: a coherent
superposition is a \emph{pure} state of zero entropy. It is therefore not the multiplicity
of possibilities that produces disorder, but the loss of coherence between them. The social
counterpart is more accurate this way: a population where everyone keeps several opinions
open is not disordered --- disorder arises from contact with an environment that fixes
positions.

\section{Thermodynamic dynamics}

\subsection{The Ising model and social temperature}

Each individual carries a binary opinion $s_i \in \{-1,+1\}$ and the population follows the
Gibbs-Boltzmann distribution \citep{ising1925}:
\begin{equation}
P(\sigma) = \frac{1}{Z}
\exp\!\left(\frac{J \sum_{\langle i,j\rangle} s_i s_j + H \sum_i s_i}{k_B T}\right),
\end{equation}
where $J$ measures local conformity, $T$ is the \emph{social temperature} (agitation,
irrationality, openness to fluctuation) and $H$ the external media field.

The model's value lies not in the analogy but in a verifiable property: the existence of a
critical temperature. In two dimensions its exact value is known \citep{onsager1944}:
\begin{equation}
\frac{T_c}{J} = \frac{2}{\ln(1+\sqrt{2})} \approx 2.269 .
\label{eq:onsager}
\end{equation}

This is the only point in this work where an exact theoretical prediction, independent of
our sociological assumptions, can validate the implementation.
Figure~\ref{fig:transition} recovers it numerically; the residual gap is a finite-size
effect on a $24\times24$ lattice.

\begin{figure}[t]
\centering
\includegraphics[width=\linewidth]{fig02_transition.png}
\caption{Consensus and susceptibility as functions of social temperature (Monte Carlo,
$24\times 24$ lattice). The susceptibility peak locates the transition. Its upward shift
relative to Onsager's value~\eqref{eq:onsager} is a finite-size effect, not an
implementation error. Figure captions are in English; notebook prose is in French.}
\label{fig:transition}
\end{figure}

\subsection{Free energy}

The system minimises its Helmholtz free energy $F = E - TS$, where $E$ represents the
tension of disagreement and $TS$ the dispersive force. As $N$ grows, configurational
entropy grows with it and the $TS$ term eventually dominates $E$.

\subsection{Social hysteresis}

This is the most directly actionable result. Below $T_c$, switching off the media field
($H \to 0$) does not return magnetisation to zero: the interaction term $-J\sum s_i s_j$
takes over and the group sustains the belief through internal cohesion.

The cycle reads in three stages, which are those of a media crisis: injection of a massive
field, official retraction, then the need for a counter-field $-H$ exceeding the group's
coercive field. Figure~\ref{fig:hysteresis} shows the cycle area decreasing with social
temperature and vanishing above $T_c$.

\begin{figure}[t]
\centering
\includegraphics[width=\linewidth]{fig05_hysteresis.png}
\caption{Left: hysteresis cycles below and above $T_c$; in the former, opinion remains
magnetised while the field is zero. Right: decay of belief memory with social temperature
--- the basis for the thermal-noise recommendation.}
\label{fig:hysteresis}
\end{figure}

\emph{Practical consequence:} passive debunking does not work. Restoring informational
neutrality leaves the belief in place; only active counter-information, or a rise in social
temperature, dissipates it.

\section{Size effects: consensus time}

The voter model \citep{clifford1973,holley1975} isolates pure social contagion, without
temperature: at each interaction an individual adopts a randomly chosen neighbour's
opinion. It yields exact scaling laws for consensus time, measured in sweeps ($N$
elementary interactions).

Measurements (figure~\ref{fig:consensus}) give an exponent $\approx 1$ in mean field and
$\approx 2$ on a ring. \textbf{A densely connected network therefore converges faster than
a geographic neighbourhood.}

This invalidates a frequent argument, that the small-world structure of social networks
\citep{watts1998} prevents consensus. Connectivity does not fragment: it accelerates and
amplifies, in whatever direction the field gives it. The causes of fragmentation are the
\emph{directional bias} of algorithmic micro-fields and the \emph{homophily} that
compartments the graph \citep{pariser2011}. The regulatory consequence is clear: the useful
lever acts on the field, not on the number of links.

\begin{figure}[t]
\centering
\includegraphics[width=\linewidth]{fig03_consensus.png}
\caption{Left: consensus time versus population size, log-log, for two topologies. Right:
trajectories of the field-favouring fraction for three bias intensities --- it is bias, not
connectivity, that decides the outcome.}
\label{fig:consensus}
\end{figure}

\section{The landscape of public opinion}

\subsection{The equation and its drift}

At macroscopic scale, collective opinion $x \in [-1,1]$ obeys a Fokker-Planck equation
\citep{risken1989}:
\begin{equation}
\frac{\partial P(x,t)}{\partial t}
= -\frac{\partial}{\partial x}\big[A(x)P(x,t)\big]
+ \frac{\partial^2}{\partial x^2}\big[B(x)P(x,t)\big].
\end{equation}

The quadratic potential $V(x) = -\frac{J}{2}x^2 - Hx$, which comes naturally to mind, can
produce \emph{no} phase transition: the corresponding exponential is convex, hence maximal
at the extremes at any temperature. The missing term is the entropy of mixing --- the number
of individual configurations compatible with a mean opinion $x$. Adding it recovers the
Helmholtz free energy:
\begin{equation}
f(x) = \underbrace{-\frac{J}{2}x^2 - Hx}_{E}
+ T\underbrace{\left[\frac{1+x}{2}\ln\frac{1+x}{2}
+ \frac{1-x}{2}\ln\frac{1-x}{2}\right]}_{-S},
\label{eq:free}
\end{equation}
giving the drift
\begin{equation}
A(x) = -f'(x) = Jx + H - T\operatorname{artanh}(x).
\label{eq:drift}
\end{equation}
The entropic restoring term bounds the dynamics within $(-1,1)$ and yields a mean-field
critical temperature $T_c = J$. The form $A(x) = Jx + H$ is its linearisation near $x = 0$.

\subsection{Diffusion and the effect of $N$}

The diffusion term reads
\begin{equation}
B(x) = \frac{k_B T (1-x^2)}{N},
\label{eq:diffusion}
\end{equation}
the factor $(1-x^2)$ switching off noise at unanimity.

The $1/N$ factor requires careful reading, since it appears to contradict the initial
hypothesis: the larger the population, the \emph{less} noisy its macroscopic variable.
There is no contradiction, but two distinct scales. Total configurational entropy is
extensive and grows as $N$; fluctuations of the mean $x$ decay as $1/N$. The defensible
thesis is therefore:

\begin{quote}
A large population does not become noisy, it becomes \emph{rigid}. Its size does not
agitate it; it deprives it of stochastic plasticity --- and it is that rigidity which makes
polarisation irreversible, since a system without noise can no longer leave the potential
well it has fallen into.
\end{quote}

This reformulation is stronger than the initial hypothesis: it explains irreversibility,
which the "entropy pump" metaphor did not.

\subsection{Three regimes}

The equilibrium distribution takes the large-deviation form
$P_{\mathrm{stat}}(x) \propto \exp(-N f(x)/T)$, where the factor $N$ makes the penalty more
severe the larger the population. Three regimes follow, obtained by varying two parameters
(figure~\ref{fig:landscape}): a fluid society ($T > T_c$), a polarised one ($T < T_c$,
$H = 0$) and a manipulated one ($T < T_c$, $H > 0$). In the last case the peak on the
field's side overwhelms the other: the "false consensus" is not an added assumption, it is
the tilted well.

The density of moderates decays exponentially with $N$. In a large population below $T_c$,
moderation is not a minority position: it is statistically inaccessible.

\begin{figure}[t]
\centering
\includegraphics[width=\linewidth]{fig04_paysage.png}
\caption{Free-energy landscape~\eqref{eq:free}, the associated stationary distributions, and
the temporal splitting of an initially moderate society. The solver uses finite volumes with
zero flux at the boundaries: probability mass is conserved to $10^{-15}$, which distinguishes
a genuine collapse of moderation from numerical leakage.}
\label{fig:landscape}
\end{figure}

\subsection{Barrier crossing}

Moving from one bloc to the other without passing through moderation is thermally activated
barrier crossing, with rate $\propto e^{-\Delta V / k_B T}$ \citep{kramers1940}, and not
tunnelling --- which has no classical counterpart. The distinction is not cosmetic:
Kramers' law supplies a temperature dependence, hence a testable prediction.

\section{Algorithmic resonance kinetics}

A piece of false information and a recommender algorithm form a positive feedback loop
analogous to the Larsen effect. Visibility $V(t)$ obeys
\begin{equation}
\ddot V + \big(\lambda - \gamma\alpha\,\sigma(V)\big)\dot V + \omega_0^2 V = \xi(t),
\label{eq:resonance}
\end{equation}
where $\lambda$ is natural damping (forgetting), $\gamma$ the algorithmic gain, $\alpha$ the
content's emotional charge, and $\sigma(V) = 1/(1+(V/V_{\mathrm{sat}})^2)$ a saturation
factor expressing the finiteness of available attention.

The instability criterion is
\begin{equation}
\boxed{\;\gamma\alpha > \lambda\;}
\label{eq:criterion}
\end{equation}
Beyond it, effective damping becomes negative: the system accumulates energy instead of
dissipating it. Figure~\ref{fig:resonance} confirms that the measured threshold coincides
with $\gamma^\star = \lambda/\alpha$.

Without saturation the unstable solution diverges as $e^{(\gamma\alpha-\lambda)t}$, which is
exact but unusable: infinite visibility does not exist and no threshold is calibratable on a
divergent model. With saturation the system settles into a limit cycle --- a self-sustaining
media oscillation of finite amplitude, matching the observed behaviour of established
disinformation.

\begin{figure}[t]
\centering
\includegraphics[width=\linewidth]{fig06_resonance.png}
\caption{Left: three fates for the same content --- damped, bounded resonant, unbounded
resonant. Right: limit-cycle amplitude versus algorithmic gain; the measured threshold
coincides with $\gamma^\star = \lambda/\alpha$.}
\label{fig:resonance}
\end{figure}

One point deserves emphasis for regulation. At \emph{uniform} gain, a more emotional item
crosses the threshold where a factual one does not. A platform therefore need not favour
disinformation in order to amplify it selectively: the bias is mechanical. This is what
makes measuring $\gamma$ more relevant, and more enforceable, than auditing editorial
intent.


\section{Empirical calibration of the amplification ratio}
\label{sec:calibration}

Criterion \eqref{eq:criterion} long remained a formal inequality: nothing indicated where
real systems sit. We estimate it here on public attention series.

\subsection{Identification}

For an isolated, non-oscillatory attention spike, \eqref{eq:resonance} reduces to its
first-order form $\dot V = \gamma\alpha\,\sigma(V)V - \lambda V$, which separates into two
measurable regimes: during the \emph{rise}, saturation is inactive and
$V \propto e^{(\gamma\alpha-\lambda)t}$; during the \emph{decay}, amplification is switched
off and $V \propto e^{-\lambda t}$. Two log-linear regressions therefore give
\begin{equation}
\lambda = r_{\mathrm{down}}, \qquad \gamma\alpha = r_{\mathrm{up}} + r_{\mathrm{down}},
\qquad \frac{\gamma\alpha}{\lambda} = 1 + \frac{r_{\mathrm{up}}}{r_{\mathrm{down}}}.
\label{eq:identification}
\end{equation}

\subsection{Data}

The series come from the Wikimedia pageviews API, filtered to the \texttt{user} agent to
exclude robots, over a \emph{pre-registered} corpus of 24 subjects split into two emotional
registers — accusation or scandal on one side, unscheduled discovery on the other. Since
Wikipedia has no recommender algorithm, the estimated $\gamma$ is an \emph{ecosystem} gain
rather than a platform's ranking function.

\subsection{Results}

Nineteen episodes are exploitable. The ratio lies between $1.5$ and $12$, with a median of
$2.5$ to $4.2$ depending on the estimator (figure~\ref{fig:calibration}). Three findings,
two of them negative.

\paragraph{The sign criterion is empty.} The ratio exceeds $1$ in every episode and under
every estimator. This is not a property of the ecosystem but a \emph{tautology of the
procedure}: an observable episode necessarily went through a growth phase, so
$r_{\mathrm{up}} > 0$. The recommendation to "prohibit configurations where
$\gamma\alpha > \lambda$" is consequently inapplicable and must be replaced by a
\textbf{ceiling} $\gamma\alpha/\lambda \leq \rho_{\max}$. This error was not detectable by
re-reading: only the attempt to measure revealed it.

\paragraph{The value is method-dependent.} With windows bounded by the return to baseline,
their length ranges from 6 to 46 days; since attention decays more slowly than an
exponential, a fit over a long window yields a smaller $\lambda$. The rank correlation
reaches $-0.94$. A second, fixed-horizon estimator makes $\lambda$ comparable across
episodes; the median nonetheless varies by a factor of $1.7$ across estimators, which forbids
anchoring a regulatory threshold to a single value.

\paragraph{The emotional-charge mechanism is not supported.} No detectable difference between
the two registers ($p \geq 0.13$), and the point estimate goes the opposite way to the
prediction. Sample sizes are small and the selection bias described below excludes the most
charged cases, so it would be excessive to call this a refutation — but the mechanism cannot
be presented as data-supported.

\paragraph{The method is blind to installed regimes.} Eleven of twenty-four subjects yield no
exploitable episode, and these are the archetypal cases: QAnon, health disinformation, vaccine
hesitancy. Their attention does not form a spike but a \emph{durable plateau}, which the
rolling baseline absorbs — prominence-based detection does not see them. The procedure
therefore selects against the very phenomenon the theory describes, which is the heaviest
limitation of this calibration.

\begin{figure}[t]
\centering
\includegraphics[width=\linewidth]{fig09_calibration.png}
\caption{An attention episode and its two exponential regimes; the ratio distribution by
emotional register; sensitivity of the median to the estimator; and the window artefact that
forced the fixed-horizon estimator.}
\label{fig:calibration}
\end{figure}



\section{The founding analogy, and its verdict}
\label{sec:analogy}

This work started from an analogy: the larger a quantum system, the faster it decoheres; the
larger a population, the harder agreement becomes. It produced measurable quantities --- a
diversity index, an amplification threshold, a critical temperature --- and that is its merit.
Its own claims were checked one by one, and none held.

\begin{center}
\begin{tabular}{@{}p{0.46\linewidth}p{0.48\linewidth}@{}}
\toprule
\textbf{What the analogy claimed} & \textbf{Verdict} \\
\midrule
The larger the system, the more disordered & \textbf{false, and backwards}: the retained
diffusion term, $B(x) = k_B T (1-x^2)/N$, makes the macroscopic variable \emph{more}
deterministic at large $N$. Two objects were conflated --- configuration entropy, extensive, and
the noise of the macroscopic variable, in $1/N$ \\[2mm]
Decoherence increases entropy & \textbf{false}: a closed system evolves unitarily. It is the
\emph{reduced} subsystem's entropy that grows \\[2mm]
$\tau_D \propto \tau_R / N$ & a \textbf{heuristic} scaling law, not Zurek's result, which
involves spatial separation and the thermal wavelength \\[2mm]
"Social tunnelling" & \textbf{metaphor}: the correct analogue is activated barrier crossing
\citep{kramers1940}, which is classical \\[2mm]
$1/k^N$ measures the improbability of consensus & describes a randomly drawn \textbf{initial
state}, not a dynamics \\
\bottomrule
\end{tabular}
\end{center}

To which is added an error that is not physical: the original account attributed fragmentation
to the connectivity of "small-world" networks \citep{watts1998}. In mean field the Voter Model
converges in $O(N)$: connectivity \emph{accelerates} consensus. Fragmentation comes from the
bias $H_i$ and from homophily, not from topology alone.

What survives --- free-energy landscape, phase transition, hysteresis, barrier crossing ---
belongs to classical statistical mechanics, in the Ising and sociophysics lineage
\citep{galam2004,castellano2009}. Nothing specifically quantum came through.

The analogy therefore worked as \textbf{scaffolding}: it suggested measurable quantities that
outlive it. That is a real role, and it is not the same as being true. Calling it "fruitful
hence sound" would be dishonest: it suggested measurements, and none of its claims survived
verification. The index was renamed accordingly --- it was the \emph{Entropic Dissipation
Index}, after decoherence, and is now the \emph{Exposed Diversity Index}, after what it
measures.

\part{Caveats}

\section{Limitations}
\label{sec:limits}

These are stated here rather than left to reviewers.

\subsection{On the index and the standard}

\paragraph{The retained form has never been measured on a real feed.} It is defined
\eqref{eq:edi}, its cost is quantified by exhaustive enumeration, and that is all. The gap is
one of \emph{data}, not of method (\S\ref{sec:logs}), and \S\ref{sec:request} states what would
close it.

\paragraph{The floor level cannot be deduced from the measurement.} This work establishes the
form of the standard and its price, not its value. Setting $0.60$ rather than $0.40$ is a
political decision no computation presented here settles.

\paragraph{Discretisation into viewpoints is a political choice.} Whoever defines the modalities
defines the index. Measuring by divergence to a declared catalogue makes the choice explicit
rather than burying it; it does not remove it.

\paragraph{The index remains gameable beyond what any automatic measure detects.} The attacks of
\S\ref{sec:adversarial} are closed, but a platform can still serve formally divergent,
substantively empty items --- bin diversity without argument diversity. Any mandated metric
becomes a target.

\paragraph{A floor is a constraint on what citizens see.} Defensible, but an intervention in
public debate, and presenting it as a neutral technical measure would be dishonest. Measuring
individual feeds also raises a privacy question that \S\ref{sec:request} answers only in part,
by aggregation.

\subsection{On the exposure measurements}

\paragraph{The form $e(R) = R^{-\eta}$ is posited.} Only its severity is estimated. A cascade
examination model, or one depending on the item, would produce a different rank dependence ---
which test \eqref{eq:exchangeability} would detect just as well, since it only tests
independence, but which the minimum detectable severity would no longer summarise.

\paragraph{Exhaustive enumeration bounds feed size.} Eight positions over four viewpoints.
Nothing guarantees that burial behaves the same on a fifty-item feed, and asserting it would
require an approximate optimisation that would itself have to be shown not to invent the result.

\paragraph{The measured datasets are slices.} A slice of Baidu-ULTR is not Baidu-ULTR; the Open
Bandit Dataset measures a fashion recommendation on three horizontal thumbnails, and
transporting its figures to a news feed is precisely the operation \S\ref{sec:logs} warns
against.

\paragraph{The ground-truth confrontation concerns a uniform target policy}, the only one this
dataset has a bucket for. Nothing guarantees importance weighting behaves as well for a target
further from the logging policy --- on the contrary, the effective sample size would be smaller
still.



\subsection{The two deepest assumptions, tested}
\label{sec:blindspots}

Two assumptions run through this entire note without ever having been tested: that examination
follows a power law, and that relevance is known. Neither is.

\paragraph{The exchangeability test survives a change of click model.} Under a \emph{cascade}
model \citep{craswell2008} --- the reader moves down the feed, stops on finding what they wanted, and otherwise
continues with probability $\gamma$ --- test \eqref{eq:exchangeability} rejects between $z = -208$
and $z = -241$ depending on $\gamma$, i.e. \emph{more strongly} than on Baidu-ULTR. That is the
result required: the test assumes nothing about the shape of attention, it tests only
independence, and its negative verdict on MIND therefore holds under any model.

\paragraph{The power law does not survive.} Under cascade, examination decays \emph{geometrically}
rather than polynomially. Simulation gives here the ground truth real data never provides --- the
positions actually examined --- and the comparison is unequivocal: at rank twelve, the fitted
$R^{-\hat\eta}$ overstates real exposure by a factor of $40$ at $\gamma = 0.95$, $129$ at
$\gamma = 0.85$ and $4{,}110$ at $\gamma = 0.60$. The fit nonetheless looks excellent, with a
standard error of $0.04$.

The consequence bears directly on \S\ref{sec:logs}: the severity $\hat\eta = 1.10 \pm 0.09$
measured on Baidu-ULTR means something \emph{only if} examination there follows a power law ---
and Baidu-ULTR is a search results page, the cascade model's home ground. Severity must therefore
be published \textbf{with its assumed shape}, and that shape remains to be tested. The
exchangeability test says whether the order carries information; it does not say in what form.

\paragraph{And estimated relevance erases the tuned methods' advantage.} The comparison in
\S\ref{sec:baselines} assumes relevance known. Having the methods rank on \emph{perceived}
relevance and judging them on the true one leaves their order unchanged --- filter and MMR remain
ahead --- but their advantage over random selection falls from $16.4$ points at zero noise to
$5.8$ points at a rank correlation of $0.50$, and \emph{inverts} below $0.36$. The reason is
mechanical: random selection is the only re-ranker that does not use relevance, hence the only one
whose shortfall is invariant to noise.

At realistic relevance, the gap between re-rankers therefore matters far less than the quality of
the underlying relevance engine --- which reinforces, by a third route, the conclusion that this
work's contribution is not its algorithm.


\subsection{A third check: the shape of examination}
\label{sec:formtest}

The exchangeability test says whether a log's order carries information; it does not say in what
form, and \S\ref{sec:blindspots} showed the form decides everything. A third check is therefore
required, and it rests on a \emph{conditional independence} rather than on a shape: under a
position model, examination of rank $R$ depends only on $R$, so the click is independent of what
happened above; under cascade, a click above suppresses examination below.

For each cell $s = (\text{item}, \text{rank})$, the $n_s$ impressions split into $n_{1s}$ preceded
by a click and $n_{0s}$ not, for $k_s$ clicks in total. Under independence those clicks distribute
as a draw without replacement:
\begin{equation}
\mathbb{E}[a_s] = \frac{k_s n_{1s}}{n_s}, \qquad
\mathbb{V}[a_s] = \frac{k_s (n_s - k_s)\, n_{1s} n_{0s}}{n_s^2 (n_s - 1)} .
\label{eq:form}
\end{equation}
This is the Mantel-Haenszel statistic, stratified by cell; conditioning removes item quality.

\paragraph{The check works.} It never rejects under a position model ($|z| < 1$ for
$\eta \in \{0, 0.5, 1, 2\}$) and rejects massively under cascade ($z = -97$ to $-307$).

\paragraph{But it does not settle what it was meant to settle.} A reader who stops clicking once
served --- while still scanning the feed --- produces the same signature: $z = -144$ under a
\emph{pure} position model with a one-click budget, against $-179$ under cascade. Nothing in the
clicks distinguishes them, as expected: in both cases the log shows that after a click there is
nothing. The difference is nonetheless decisive, since under a budget the item placed below a
click is examined, and under cascade it is not.

\paragraph{On Baidu-ULTR, no cascade signature.} The standardised deviation there is
\emph{positive} --- $+5.6$ at the most permissive threshold, $+0.5$ to $+0.6$ at the strictest ---
whereas a cascade would make it negative. The sign is that of the heterogeneity confounder, which
the test announced as its limit. Three readings remain open: examination there is close to a
position model; a cascade exists but the confounder masks it; or coverage is too thin, $108$ cells
at the strictest threshold.

\paragraph{A protocol error, caught.} Restricting the log to multi-click sessions --- where a
one-click budget is excluded by construction --- looked like the obvious way to separate the two
models, and gives $z = -8.2$ ($p = 2 \times 10^{-16}$) on Baidu-ULTR. Yet a feed's click count is
a \emph{collider} of its individual clicks: conditioning on it induces negative dependence between
them, and therefore manufactures the sought signature. On a log simulated under a pure position
model, the same restriction moves $z$ from $+0.74$ to $-45.7$. What was measured was not a
property of the log but the selection performed.

\paragraph{What would settle it.} A measure of \emph{examination} rather than of clicks --- display
time, scroll depth, items slipped off screen. Baidu-ULTR publishes those columns; this work has
not read them.


\subsection{Exposure measured rather than estimated}
\label{sec:measuredexposure}

Baidu-ULTR publishes a \texttt{displayed\_time} column --- the display time of each document ---
which this work had not read. It measures directly what was \emph{shown}, where everything above
observed only what was \emph{clicked}.

\paragraph{It lifts the ambiguity of \S\ref{sec:formtest}.} On clicks, cascade and budget are
conflated; on examination, separation is total: $z = -158$ to $-412$ under cascade against
$z = -1.03$ under budget, whatever the budget. Applied to Baidu-ULTR, the measure \textbf{refutes
cascade}: examination below a click is there \emph{more} frequent, not less --- $0.221$ against
$0.136$ at rank nine, $z = +8.4$ to $+10.9$ depending on threshold. This is the heterogeneity
confounder, measured directly.

\paragraph{It revises the published figure.} Exposure severity, measured by document fixed effects
on display, is $0.882 \pm 0.046$ over \textbf{143} documents, against $1.085 \pm 0.093$ over $55$
for the click-based estimate. The estimate \textbf{overstates decay by $23\%$}, and the gap is not
imprecision but an identified confounder: a click is the product of examination and attractiveness,
and attractiveness also declines with rank.

\paragraph{It calls into question the form used throughout.} Fitted on the \emph{measured} curve,
the three candidate forms give $R^2 = 0.72$ for the power law --- with $42\%$ maximum error ---,
$0.96$ for geometric decay and $0.99$ for a \emph{screen fold} model. The $R^{-\eta}$ law is the
worst of the three, and it errs most where the standard is decided: at rank two it predicts $0.47$
where measurement gives $0.88$.

\paragraph{Consequence for \S\ref{sec:request}.} The four tables requested bear on clicks. One more
column --- the number of impressions actually \emph{displayed} per cell --- would remove at a
stroke the need to estimate $\eta$, the shape assumption that accompanies it, and the
attractiveness confounder. It is moreover less sensitive than clicks: knowing an item was displayed
says less about a reader than knowing they chose it.

\paragraph{Reservation.} A positive display time attests that the document was shown, not that it
was looked at. The measure therefore overstates exposure, by an amount this work cannot quantify.


\subsection{Format, return, and an explanation withdrawn}
\label{sec:formatreturn}

Two Baidu-ULTR columns remained unexploited: \texttt{slipoff\_count\_after\_click}, which counts
documents that slipped off screen after a click, and \texttt{media\_type}, which distinguishes
display formats.

\paragraph{Cascade refuted a second time.} The first column is non-zero for $43.7\%$ of clicked
rows and $0.5\%$ of the rest: after a click, the log explicitly records that documents scrolled
past, hence that the reader came back. Under a strict cascade it would always be zero. The share
moreover grows with the depth of the click --- $0.37$ at rank one, $0.59$ at rank six. The
refutation of \S\ref{sec:measuredexposure}, obtained by a statistical test, is thus confirmed by a
recorded fact, which is an independent route.

\paragraph{Format alters exposure, and not through geometry.} A rich format --- answer box, image,
table --- is taller: $273$ pixels median against $192$. One would therefore expect a pushing-down
effect. At equal rank \emph{and} comparable cumulative height, a rich format above nonetheless
removes $8.4$ points of examination from what follows, median over $21$ strata all of the same
sign. The effect therefore does not run through space occupied, and the remaining explanation is
\emph{satisfaction}: an answer box answers the question, and the reader stops descending.

The consequence bears on any rank-only measure, including ours: exposure at rank $R$ depends on
the \textbf{format of what is above it}, and no law $e(R)$ can represent it since it does not take
the page as an argument. The attention discount $w_R$ of \eqref{eq:distribution} is therefore not
a property of the surface but of the \textbf{page served}: two feeds of identical composition,
served with different formats, do not expose the same thing.

\paragraph{And one explanation in this note is withdrawn.} \S\ref{sec:measuredexposure} attributed
the two-stage shape of the measured curve to a \emph{screen fold}, invoking the
\texttt{serp\_height} column. That column allows the hypothesis to be tested, and it refutes it:
rank explains examination better (McFadden $R^2$ $0.082$) than cumulative height above ($0.057$),
and pixels add only $0.0007$ once rank is known. The finding --- the curve is not a power law ---
holds; the mechanical explanation falls.

\paragraph{Consequence for \S\ref{sec:request}.} The cell table already asks for a displays column;
one more is needed, the \textbf{format} of the item served. Without it, two feeds of identical
composition can expose different diversities with nothing signalling it. It is moreover the least
sensitive of the requested columns: an item's format says nothing about its reader.

\subsection{What the counter-expertise withdrew or restricted}
\label{sec:counterexpertise}

The measuring instruments were confronted with the field's literature and with four
counter-tests. Three results of this note come out modified.

\paragraph{A withdrawn conclusion.} The comparison of the four diversity measures
(\S\ref{sec:adversarial}) was carried out \emph{at the same nominal floor}. Since the measures do
not live on the same scale, that comparison is meaningless: at equal actually exposed diversity,
the two retained measures cost the same, within $0.6\%$ of the exact bound. The claim that
proximity to a target "resists" burial better is a scale artefact. What makes the standard is the
\textbf{level} required and \textbf{rank-awareness}, not the choice of measure.

\paragraph{A restricted conclusion.} The figure "certified at $0.70$, it exposes $0.36$" assumes a
$1/R$ attention discount, i.e. $\eta = 1$. Recomputed with the \emph{measured} severity of a
three-thumbnail banner ($\eta \approx 0.1$), the same feed exposes $0.747$: burial disappears. At
$\eta = 2$ it exposes only $0.157$, and the loophole costs $2.7\%$ instead of $20.8\%$. The
attention discount must therefore be \textbf{measured on the surface}, and the stable price is
that of the rank-aware standard --- $20.8$ to $24.1\%$ depending on the surface.

\paragraph{A widened uncertainty.} The estimator of \S\ref{sec:exposure} assumes a multiplicative
click model. The literature documents an \emph{affine} one, where readers click out of trust on
irrelevant top-ranked items \citep{vardasbi2020}, and proves that inverse propensity weighting
cannot correct it. Under that model the estimated severity drifts by $+12.8\%$ without the
standard error signalling it. Far less severe than positing $\eta$ blindly, but the published
interval is too narrow.

\paragraph{A confirmed recommendation.} The doubly robust estimator does \emph{worse} than plain
weighting on the Open Bandit Dataset ($-4.9\%$ against $+2.5\%$), and the effective sample size
does not move: it depends only on the importance weights. No estimator replaces exploration.

\paragraph{And a caution from the literature.} On Baidu-ULTR --- the very dataset where this note
measures $\hat\eta = 1.10$ --- \citet{hager2024} find that position-bias corrections improve click
prediction \emph{without} improving ranking quality as judged by expert annotators. The presence
of the bias is confirmed by four concordant methods, which corroborates our figure; the next step
does not follow mechanically.

\subsection{On the model}

\paragraph{Calibration barely begun.} A single composite parameter is estimated
(\S\ref{sec:calibration}); $J$, $T$, $\gamma$ and $\alpha$ taken separately have no estimation
procedure. The most promising route would be inference of Kramers-Moyal coefficients on long
opinion series, which would allow the drift to be \emph{tested} rather than fitted.

\paragraph{The model's one proper prediction is not verified.} The emotional-charge mechanism
$\alpha$ predicts a difference between registers --- amplification, persistence, switching rate.
Four independent measurements, including a replication across $440$ subjects and a blind
double-coded annotation, find no trace of it. The apparent gap in the pilot corpus was a
selection artefact.

\paragraph{An analogy is not an explanation.} Nothing here demonstrates that opinions
\emph{obey} statistical mechanics; the work establishes that a formalism borrowed from physics
\emph{reproduces} some of its behaviours. On the founding analogy itself, the verdict is
sharper still (\S\ref{sec:analogy}).

\paragraph{Individuals are not spins.} Intentionality, irony and strategic anticonformism have
no equivalent in a spin model; opinions are multidimensional \citep{deffuant2000}. Above all,
the environment is not passive: an algorithm optimises a cost function, which makes it a
strategic actor rather than a thermal bath. That is the analogy's weakest point --- and it is
also what makes the filter of \S\ref{sec:filter} conceivable: a cost function can be changed.

\paragraph{Scope of the simulations.} The regimes explored were chosen for legibility. No
systematic sensitivity study was carried out, system sizes are modest, and the conclusions of
Part~\ref{part:model} are qualitative: they concern the existence of regimes and the direction
of dependencies, never transferable numerical values.

\subsection{On the annotation}

\paragraph{The coders are not independent.} Blind double coding of the corpus yields a Fleiss
$\kappa$ of $0.921$, but the three coders are instances of the same language model: their
agreement necessarily overstates what independent human judges would produce. It is the last
methodological reservation that no computation lifts.

\section{Reproducibility}

The entire work is openly available and runs in containers, with $556$ automated tests. The
structural checks cover Onsager's critical temperature \eqref{eq:onsager}, exact conservation of
probability mass, the scaling exponents of consensus time, the strict positivity of the
hysteresis loop area below $T_c$, the equivalence of filter \eqref{eq:score} with an engagement
filter at $\mu = 0$, the sign of the exchangeability test's rejection on a genuinely ordered
log, and the exactness of the measurements recomputed from the aggregate tables of
\S\ref{sec:request}.

The impression logs used are not redistributed --- they weigh $135$ MB to $2.2$ GB and each
carries its own licence --- but the repository keeps verified digests of them, which reproduce
every figure published here identically. Every figure is regenerated by a notebook.

\begin{center}
\url{https://github.com/s-geffroy/Indice-Diversite-Exposee}\\
\url{https://s-geffroy.github.io/Indice-Diversite-Exposee/}
\end{center}

This work is open to critical review. The most useful feedback would concern the retained form
of the index and the level of any floor, or the measurement instruments of \S\ref{sec:exposure}
and \S\ref{sec:evaluation}, where a methodological error would be the most costly.

\bibliography{refs}

\end{document}
